\chapter{Public Health and Preventive Medicine}

\medterm{Adequate intake}-- An estimate of the amount of a nutrient needed by healthy people. The Adequate Intake is used when there isn’t enough information to set a recommended dietary allowance (RDA).

\medterm{Antimicrobial Resistance}-- Development of resistance to antimicrobial agents. Occurs when microorganisms evolve mechanisms to survive exposure to drugs that previously killed or inhibited them. Driven by overuse and misuse of antibiotics in humans, agriculture, and animal husbandry. Results in treatment failure, prolonged illness, increased healthcare costs, and higher mortality. Major global health threat requiring antimicrobial stewardship, surveillance, and development of new agents.

\medterm{Asbestosis}-- Occupational lung disease caused by inhalation of asbestos fibers. Characterized by diffuse interstitial pulmonary fibrosis. Long latency period of 10 to 20 years. Associated with increased risk of mesothelioma and lung cancer.

\medterm{Booster Dose}-- Additional dose of a vaccine given after the primary series to re-stimulate the immune response and maintain protection. Required for vaccines where immunity wanes over time such as tetanus and pertussis.

\medterm{Burden of Disease}-- Measure of the gap between current health status and an ideal health situation where everyone lives to old age free of disease and disability. Measured in DALYs. Used to set health priorities and allocate resources.

\medterm{Burden of Illness}-- Total impact of a disease on a population measured in terms of mortality, morbidity, disability, and economic costs. Used to set health priorities and allocate research funding.

\medterm{Byssinosis}-- Occupational lung disease caused by inhalation of cotton, flax, or hemp dust. Characterized by chest tightness and wheezing that improves away from work. Also known as brown lung disease.

\medterm{Carbon Monoxide Poisoning}-- Toxic effect of inhaling carbon monoxide an odorless, colorless gas. Binds hemoglobin with high affinity causing tissue hypoxia. Common sources include faulty furnaces, car exhaust, and generators.

\medterm{Contraindication to Vaccination}-- Condition or factor that increases the risk of a serious adverse reaction to a vaccine. Absolute contraindications include severe allergic reaction to a previous dose. Relative contraindications include immunosuppression for live vaccines.

\medterm{Crossover Trial}-- Study design where participants receive both treatments in a random sequence with a washout period between. Each participant serves as their own control. Increases statistical power but requires a washout period to avoid carryover effects.

\medterm{Disease Prevention}-- Actions taken to prevent disease or detect it early before symptoms develop. Includes immunization, screening, and chemoprophylaxis. Distinguished from treatment which occurs after disease has manifested.

\medterm{Disease Surveillance}-- Continuous systematic collection, analysis, interpretation, and dissemination of health data for disease prevention and control. Essential for early detection of outbreaks and monitoring disease trends.

\medterm{Effect Modification}-- Situation where the magnitude of the association between exposure and outcome differs across levels of another variable. Also called interaction. Unlike confounding, effect modification is a real biological phenomenon that should be reported.

\medterm{Effect Size}-- Quantitative measure of the magnitude of an observed effect. Common measures include Cohen’s d for mean differences and odds ratios for binary outcomes. Important for assessing clinical significance beyond statistical significance.

\medterm{Epidemic Curve}-- Graphical representation of the number of cases of disease over time during an outbreak. Can be point source, propagated, or intermittent. Provides information about the outbreak pattern, incubation period, and potential source.

\medterm{Epidemics}-- The occurrence of cases of a disease or health-related event in a community or region at a rate clearly in excess of normal expectancy. Epidemics arise from point sources (single exposure, e.g., contaminated food at an event), propagated spread (person-to-person transmission, e.g., influenza, measles), or mixed patterns. Investigation steps: confirm diagnosis, establish case definition, verify cases, plot epidemic curve (cases over time), calculate attack rates, generate hypotheses about source and mode of transmission, implement control measures (isolation, quarantine, vaccination, treatment, sanitation), and communicate findings. A pandemic is an epidemic that has spread across multiple continents. Epidemiology is the scientific discipline that studies the distribution and determinants of health-related states in populations.

\medterm{Epidemiological study}-- An investigation of the links between certain behaviors or risk factors and the occurrence of disease or good health in a population.

\medterm{Epidemiologic Transition}-- Theory describing the shift in patterns of mortality and disease from infectious to chronic noncommunicable diseases as countries develop. Characterized by changing age distributions and life expectancies.

\medterm{Essential Fatty Acid}-- Fatty acids that cannot be synthesized by the body and must be obtained from the diet. Include omega-3 and omega-6 fatty acids. Important for brain development, inflammation regulation, and cardiovascular health.

\medterm{Folate Deficiency}-- Causes megaloblastic anemia and neural tube defects in developing fetuses. Prevention through folic acid supplementation before and during pregnancy. Found in leafy greens, legumes, and fortified grains.

\medterm{Genetic Epidemiology}-- Study of the role of genetic factors in determining health and disease in populations. Includes twin studies, family studies, and genome-wide association studies.

\medterm{Geriatrician}-- A physician who specializes in the care of older patients.

\medterm{Hazard Ratio}-- Ratio of hazard rates between two groups at any point in time. Used in survival analysis to compare the rate of events between treatment and control groups. Similar to relative risk but accounts for time-to-event data.

\medterm{Health Equity}-- Attainment of the highest level of health for all people. Involves addressing social determinants such as income, education, housing, and access to healthcare that produce systematic disparities between population groups. Requires targeted policies and interventions to eliminate avoidable and unfair differences in health outcomes. Distinct from health equality in that equity allocates resources according to need.

\medterm{Health Inequality}-- Differences in health status between different population groups that are avoidable and unfair. Includes inequalities in access to healthcare, healthy food, safe housing, and education.

\medterm{Herd Immunity}-- Protection of a population against an infectious disease when a sufficient proportion is immune either through vaccination or prior infection. The threshold varies by disease. Reduces the risk of infection for susceptible individuals.

\medterm{Inactivated Vaccine}-- Vaccine containing killed pathogens that cannot replicate. Safer than live-attenuated vaccines but require multiple doses and booster shots to maintain protection. Induce humoral immune responses primarily. Examples: influenza (injectable), polio (IPV), hepatitis A, rabies. More stable for storage and transport than live vaccines.

\medterm{Iodine Deficiency}-- Leading cause of preventable intellectual disability worldwide. Causes goiter, hypothyroidism, and cretinism in severe cases. Prevented through iodised salt supplementation programs. Universal salt iodisation is the recommended strategy endorsed by WHO and UNICEF. Pregnant women and young children are most vulnerable to the effects of deficiency.

\medterm{Iron Deficiency}-- Most common nutritional deficiency worldwide. Causes microcytic hypochromic anemia, fatigue, and impaired cognitive development. Risk factors include menstruation, pregnancy, and inadequate dietary intake.

\medterm{Isolation}-- Separation of individuals with a confirmed communicable disease from healthy individuals. Aims to prevent transmission to others. Required for diseases such as tuberculosis, measles, and COVID-19.

\medterm{Lead Exposure}-- Environmental health hazard from lead in paint, water, soil, and dust. Particularly harmful to children causing developmental delays and behavioral problems. Primary prevention through lead abatement and water pipe replacement.

\medterm{Lead Poisoning}-- Occupational and environmental hazard from exposure to lead. Affects the nervous system, kidneys, and hematopoietic system. Particularly dangerous for children causing developmental delays. Measured by blood lead levels.

\medterm{Lead Time}-- Period between detection of disease through screening and the time it would have been diagnosed clinically. Longer lead time does not necessarily mean better prognosis. Used to calculate apparent survival benefit from screening.

\medterm{Live-Attenuated Vaccine}-- Vaccine containing weakened forms of the pathogen that can still replicate but do not cause disease in immunocompetent individuals. Provides strong and long-lasting immunity. Contraindicated in immunocompromised patients.

\medterm{Low-calorie diet}-- A weight-loss plan that limits calories to 800–1,500 a day.

\medterm{Mercury Exposure}-- Environmental health hazard from mercury in fish, industrial emissions, and dental amalgams.

\medterm{Mesothelioma}-- Rare cancer of the mesothelial lining strongly associated with asbestos exposure. Latency period of 20 to 40 years. Poor prognosis with median survival less than one year. Mostly affects the pleura.

\medterm{Meta-Analysis}-- A statistical technique that combines and analyzes the results of multiple independent studies on the same topic to produce a single, more precise estimate of the treatment effect or association. Meta-analysis increases statistical power, improves precision, resolves uncertainty from conflicting studies, and can identify subgroup effects. Key components: systematic literature search, quality assessment of included studies, extraction of effect estimates, calculation of a weighted pooled effect (fixed-effect or random-effects model), assessment of heterogeneity (I² statistic, Cochrane Q test), and publication bias assessment (funnel plot, Egger's test). Forest plot: graphical display showing individual study results and the pooled estimate. Limitations: garbage in, garbage out (quality depends on quality of included studies), publication bias (positive results more likely to be published), and heterogeneity (clinical and methodological diversity). Meta-analysis is a key tool in evidence-based medicine and forms the highest level of evidence (systematic reviews/meta-analyses on top of the evidence pyramid).

\medterm{Multimorbidity}-- Co-occurrence of two or more chronic conditions in an individual. Increasingly common with aging populations. Affects quality of life, healthcare utilization, and treatment complexity.

\medterm{Natural Experiment}-- Observational study in which the exposure of interest is not assigned by the researcher but occurs naturally due to external factors such as policy changes, natural disasters, or geographical variation. The researcher leverages this natural variation to estimate causal effects. Strength: allows study of exposures that cannot be randomised for ethical or practical reasons. Limitation: lacks the randomisation of controlled trials, making confounding and bias more difficult to control. Techniques such as difference-in-differences, instrumental variables, and regression discontinuity are used to strengthen causal inference.

\medterm{Negative Predictive Value}-- Proportion of negative test results that are true negatives. Depends on disease prevalence in the tested population. Higher prevalence decreases NPV. Important for ruling out disease.

\medterm{Neglected Tropical Diseases}-- Group of communicable diseases that prevail in tropical and subtropical conditions affecting the poorest populations. Include Chagas disease, dengue, leishmaniasis, and schistosomiasis. Targeted for control and elimination.

\medterm{Neyman Bias}-- Type of selection bias where cases with rapid onset or fatal outcome are missed because they occur and resolve between screening intervals. Also called prevalence-incidence bias. Affects estimates of disease severity from screening studies.

\medterm{Noise-Induced Hearing Loss}-- Permanent sensorineural hearing loss caused by prolonged exposure to loud noise. Most common occupational disease worldwide. Preventable through engineering controls, administrative controls, and hearing protection.

\medterm{Notifiable Diseases}-- Diseases required by law to be reported to public health authorities. Reporting enables tracking of disease trends and outbreak detection. Examples include tuberculosis, measles, and sexually transmitted infections, foodborne illnesses. Reporting enables public health surveillance, outbreak detection, and implementation of control measures. Lists vary by jurisdiction.

\medterm{Occupational Carcinogen}-- Agent in the workplace that increases the risk of cancer. IARC classifies carcinogens into groups from 1 (definitely carcinogenic) to 3 (not classifiable). Includes asbestos, benzene, chromium VI, and formaldehyde.

\medterm{Occupational Health}-- Branch of public health concerned with the health and safety of workers. Includes prevention and management of workplace injuries and illnesses, ergonomic assessment, and promotion of worker wellbeing.

\medterm{Odds Ratio}-- Ratio of the odds of exposure among cases to the odds of exposure among controls. Approximates relative risk when disease is rare. Widely used in case-control studies where direct incidence cannot be calculated.

\medterm{Opportunity Cost}-- Value of the next best alternative foregone when making a choice. In healthcare, represents the benefits that could have been gained from investing resources in a different intervention. Important concept in health economics. When choosing between healthcare interventions, the opportunity cost is the health benefits foregone from not funding the next best alternative. Used in cost-effectiveness analysis and health technology assessment to inform resource allocation decisions.

\medterm{Outbreak Investigation}-- Systematic process to identify the source, mode of transmission, and factors contributing to an epidemic. Aims to implement control measures and prevent future outbreaks. Follows a standard series of steps.

\medterm{Overdiagnosis}-- Detection of disease through screening that would never have caused symptoms or death during the patient’s lifetime. Results in unnecessary treatment and associated harms. Particularly problematic in cancer screening.

\medterm{Overtreatment}-- Treatment of conditions detected through screening that would not have caused harm if left undetected.

\medterm{Pandemic Preparedness}-- Planning and capacity building to respond to pandemic threats. Includes surveillance, laboratory capacity, vaccine development, and healthcare system resilience. Guided by the International Health Regulations (IHR 2005). Includes early warning systems, stockpiling of vaccines and antivirals, surge capacity planning, and coordination mechanisms. Lessons from COVID-19, influenza pandemics, and emerging infectious disease threats inform ongoing preparedness efforts.

\medterm{Patient-Centered Care}-- Approach to healthcare that is respectful of and responsive to individual patient preferences, needs, and values. Ensures patient values guide all clinical decisions. Includes shared decision-making, communication, and care coordination. Improves patient satisfaction, treatment adherence, and health outcomes.

\medterm{Peer Review}-- The process by which research submitted for publication or presentation is evaluated by independent experts (peers) in the same field to assess its quality, validity, originality, and significance before acceptance. Types: single-blind (reviewers know authors, but not vice versa), double-blind (neither knows the other), open (both identities are known), and post-publication review. In healthcare, peer review also applies to clinical practice: regular review of doctors' performance by colleagues to maintain and improve standards of care (including case note review, morbidity and mortality meetings, and clinical audit). Peer review is the cornerstone of scientific quality control but is not infallible — it can be subject to bias, delay, and failure to detect major errors or fraud.

\medterm{Personal Protective Equipment}-- Equipment worn to minimize exposure to workplace hazards. Includes respirators, gloves, eye protection, and hearing protection. Last line of defense in the hierarchy of controls.

\medterm{Primary Prevention}-- Prevention of disease before it occurs. Includes health promotion and specific protection measures such as vaccination, health education, and environmental modifications. Aims to reduce the incidence of disease.

\medterm{Primordial Prevention}-- Prevention of the development of risk factors in populations through environmental, economic, social, and behavioural interventions before risk factors emerge. Targets upstream determinants such as poverty, education, and urban planning. Example: policies to reduce air pollution to prevent cardiovascular and respiratory disease. Distinct from primary prevention which targets established risk factors.

\medterm{Propagated Outbreak}-- Outbreak where disease spreads from person to person. Characterized by successive waves of cases corresponding to generations of transmission. The epidemic curve shows multiple peaks.

\medterm{Proportional Mortality Ratio}-- Proportion of deaths due to a specific cause relative to all deaths in a population. Useful for comparing the contribution of different causes of death.

\medterm{Prospective Cohort Study}-- Observational study following a group of individuals over time to determine who develops the outcome of interest. Measures incidence and relative risk. Strong for studying rare exposures. Expensive and time-consuming.

\medterm{Publication Bias}-- Tendency for studies with positive or significant results to be published more readily than those with negative or null results. Can distort meta-analyses and systematic reviews. Addressed through trial registration and comprehensive literature searching.

\medterm{Quarantine}-- Separation and restriction of movement of healthy individuals who may have been exposed to a communicable disease. Aims to prevent disease transmission. Applied during the incubation period to monitor for symptom development. Historically used for plague, cholera, and yellow fever. Legal authority for quarantine varies by jurisdiction. Distinguished from isolation (separation of those already confirmed ill).

\medterm{Quasi-Experimental Design}-- Study design that lacks randomization but includes an intervention. Uses comparison groups that are not randomly assigned. Susceptible to confounding but can be strengthened by techniques such as difference-in-differences analysis.

\medterm{Radon Exposure}-- Natural radioactive gas from the decay of uranium in soil and rock. Second leading cause of lung cancer after smoking. Enters buildings through cracks in foundations. Tested with radon detectors and mitigated by ventilation systems.

\medterm{Relative Risk}-- Ratio of the incidence rate of disease in the exposed population to the incidence rate in the unexposed population. A relative risk greater than 1 indicates increased risk with exposure; less than 1 indicates a protective effect. Used in cohort studies and clinical trials. Important for establishing causal associations. Also called risk ratio.

\medterm{Relative Risk Reduction}-- Proportional reduction in risk between the treatment and control groups. Calculated as the difference in risk divided by the control group risk. Often used in clinical trial reporting.

\medterm{Risk Communication}-- Process of exchanging information about health risks between experts and the public. Includes risk assessment, risk perception, and risk management communication. Important during outbreaks and public health emergencies.

\medterm{Risk Factor}-- Attribute or exposure that is associated with an increased probability of disease. Does not necessarily imply causation. Can be modifiable or nonmodifiable. Important for disease prevention and health promotion.

\medterm{Scientific Evidence}-- Information gathered through systematic observation, experimentation, and reasoning that supports or refutes a hypothesis or theory. In medicine, scientific evidence forms the basis for clinical practice and is ranked by quality (hierarchy of evidence). Levels (from highest to lowest): systematic reviews and meta-analyses, randomized controlled trials (RCTs), cohort studies, case-control studies, cross-sectional studies, case series/case reports, expert opinion, and anecdotal evidence. Evidence-based medicine integrates the best available scientific evidence with clinical expertise and patient values and preferences. The GRADE system rates the quality of evidence as high, moderate, low, or very low based on study design, risk of bias, consistency, directness, and precision.

\medterm{Sensitivity}-- Proportion of individuals with disease who test positive. True positive rate. Important for screening tests where missing cases has serious consequences.

\medterm{Shared Decision-Making}-- Process where healthcare providers and patients collaborate to make clinical decisions. Involves providing evidence-based information about options, eliciting patient preferences, and reaching a mutually agreed plan.

\medterm{Silicosis}-- Occupational lung disease caused by inhalation of crystalline silica dust. Characterized by progressive fibrosis of the lungs. Increased risk of tuberculosis and lung cancer. Common in miners, sandblasters, and stone cutters.

\medterm{Subunit Vaccine}-- Vaccine containing specific components of the pathogen such as proteins or polysaccharides. Cannot cause disease because they do not contain live organisms. Examples include hepatitis B and HPV vaccines.

\medterm{Syndromic Surveillance}-- Monitoring of health-related data that precede diagnosis and signal a potential outbreak. Includes monitoring of symptoms, school absenteeism, and over-the-counter medication sales. Provides earlier warning than traditional surveillance.

\medterm{Systematic Review}-- A structured, comprehensive review of the literature that uses explicit, reproducible methods to identify, select, appraise, and synthesize all available evidence on a specific research question. Key features: pre-specified protocol (PROSPERO registration), comprehensive search strategy (multiple databases: MEDLINE, EMBASE, Cochrane CENTRAL, CINAHL), clear inclusion/exclusion criteria, quality assessment of included studies (risk of bias tools: Cochrane RoB 2 for RCTs, ROBINS-I for non-randomized studies), data extraction, and synthesis (narrative or quantitative via meta-analysis). PRISMA (Preferred Reporting Items for Systematic Reviews and Meta-Analyses) guidelines standardize reporting. Systematic reviews sit at the top of the evidence pyramid and are essential for clinical guideline development and evidence-based practice.

\medterm{Tertiary Prevention}-- Prevention of disease progression and complications after clinical disease has been diagnosed. Includes rehabilitation, disease management, and secondary prevention of complications. Aims to reduce disability and improve quality of life.

\medterm{Vaccination}-- Administration of a vaccine to stimulate the immune system and provide protection against a specific infectious disease. Can be live-attenuated, inactivated, subunit, toxoid, or mRNA-based. Most cost-effective public health intervention.

\medterm{Vaccine Effectiveness}-- Measure of vaccine performance in real-world conditions outside clinical trials. Lower than efficacy due to factors such as vaccine storage, administration, and population heterogeneity. Important for public health decision-making.

\medterm{Vaccine Efficacy}-- Proportionate reduction in disease incidence among vaccinated individuals compared to unvaccinated individuals under ideal conditions. Measured in clinical trials. Different from effectiveness which measures real-world performance.

\medterm{Vitamin A Deficiency}-- Leading cause of preventable childhood blindness worldwide. Causes xerophthalmia, keratomalacia, and increased susceptibility to infections. Common in developing countries with limited access to animal sources and orange fruits and vegetables.

\medterm{Vitamin D Deficiency}-- Causes rickets in children and osteomalacia in adults. Risk factors include limited sun exposure, dark skin, and inadequate dietary intake. Prevented through sunlight exposure, dietary sources, and supplementation.

\medterm{Zinc Deficiency}-- Impairs immune function, growth, and wound healing. Common in developing countries. Risk factors include vegetarian diets, malabsorption, and chronic diseases. Manifests with growth retardation, diarrhea, and recurrent infections.

